\documentclass[11pt,a4paper]{article}

\usepackage[utf8]{inputenc}
\usepackage[T1]{fontenc}
\usepackage{lmodern}
\usepackage{amsmath,amssymb,amsfonts}
\usepackage{booktabs}
\usepackage{hyperref}
\usepackage[margin=1in]{geometry}
\usepackage{microtype}
\usepackage{xcolor}

\hypersetup{
  colorlinks=true,
  linkcolor=blue!60!black,
  citecolor=blue!60!black,
  urlcolor=blue!60!black,
}

\title{\textbf{Supplementary Material: The Torus Electron}}
\author{James P.\ Galasyn and Claude Th\'eodore}
\date{}

\begin{document}
\maketitle

%----------------------------------------------------------------------
\section{S1. Torus Knot Parameterization}
\label{sec:S1}

A $(p,q)$ torus knot on a torus with major radius $R$ and minor radius $r$ is parameterized by $\lambda \in [0, 2\pi)$:
\begin{align}
  x(\lambda) &= (R + r\cos q\lambda)\cos p\lambda, \\
  y(\lambda) &= (R + r\cos q\lambda)\sin p\lambda, \\
  z(\lambda) &= r\sin q\lambda.
\end{align}

The toroidal angle advances as $\theta = p\lambda$ and the poloidal angle as $\phi = q\lambda$. For $(p,q) = (2,1)$, the curve winds twice toroidally for each poloidal circuit, closing after $\lambda$ traverses $[0, 2\pi)$.

\paragraph{Tangent vectors.} Differentiating:
\begin{align}
  \frac{dx}{d\lambda} &= -rq\sin(q\lambda)\cos(p\lambda) - (R + r\cos(q\lambda))\,p\sin(p\lambda), \\
  \frac{dy}{d\lambda} &= -rq\sin(q\lambda)\sin(p\lambda) + (R + r\cos(q\lambda))\,p\cos(p\lambda), \\
  \frac{dz}{d\lambda} &= rq\cos(q\lambda).
\end{align}

\paragraph{Arc length element.} After expanding and simplifying (cross terms cancel in $dx^2 + dy^2$):
\begin{equation}
  \left|\frac{d\mathbf{r}}{d\lambda}\right|^2 = r^2 q^2 + (R + r\cos q\lambda)^2 p^2.
\end{equation}

In the thin-torus limit $r \ll R$:
\begin{equation}
  \left|\frac{d\mathbf{r}}{d\lambda}\right| \approx \sqrt{p^2 R^2 + q^2 r^2}\left(1 + \frac{Rrp^2\cos(q\lambda)}{p^2 R^2 + q^2 r^2}\right).
\end{equation}

The cosine term integrates to zero over a full period, giving the total path length:
\begin{equation}
  L = \int_0^{2\pi}\left|\frac{d\mathbf{r}}{d\lambda}\right| d\lambda \approx 2\pi\sqrt{p^2 R^2 + q^2 r^2}.
\end{equation}

For $(2,1)$ with $r = \alpha R$: $L \approx 2\pi R\sqrt{4 + \alpha^2} \approx 4\pi R$.

%----------------------------------------------------------------------
\section{S2. Self-Energy Derivation}
\label{sec:S2}

\subsection{Neumann inductance}

The magnetic self-energy of a current loop on a torus is determined by the Neumann inductance formula~\cite{jackson1999}. For a thin torus ($r \ll R$), the inductance of the toroidal current path is:
\begin{equation}
  \mathcal{L} = \mu_0 R\left[\ln\frac{8R}{r} - 2\right].
\end{equation}

This is the standard result for a circular loop of radius $R$ with wire radius $r$ (Jackson \S5.6~\cite{jackson1999}).

\subsection{Current and self-energy}

A charge $e$ traversing a closed path of length $L$ at speed $c$ creates a time-averaged current:
\begin{equation}
  I = \frac{e}{T} = \frac{ec}{L}.
\end{equation}

The magnetic self-energy is:
\begin{equation}
  U_{\text{mag}} = \frac{1}{2}\mathcal{L}I^2 = \frac{\mu_0 R}{2}\left[\ln\frac{8R}{r} - 2\right]\left(\frac{ec}{L}\right)^{\!2}.
\end{equation}

\subsection{Electric--magnetic equality at \texorpdfstring{$v = c$}{v = c}}

For a charge moving at speed $v$, the electromagnetic field in the lab frame has energy densities:
\begin{equation}
  u_B = \frac{B^2}{2\mu_0}, \qquad u_E = \frac{\varepsilon_0 E^2}{2}.
\end{equation}

In the ultrarelativistic limit $v \to c$, the field becomes purely transverse with $|\mathbf{E}| = c|\mathbf{B}|$, giving $u_E = u_B$. The total electromagnetic self-energy is therefore:
\begin{equation}
  U_{\text{EM}} = U_{\text{mag}} + U_{\text{elec}} = 2U_{\text{mag}} = \mu_0 R\left[\ln\frac{8R}{r} - 2\right]\left(\frac{ec}{L}\right)^{\!2}.
\end{equation}

\subsection{Ratio to circulation energy}

The circulation energy is $E_{\text{circ}} = 2\pi\hbar c / L$. The ratio:
\begin{equation}
  \frac{U_{\text{EM}}}{E_{\text{circ}}} = \frac{\mu_0 R[\ln(8R/r) - 2](ec)^2 / L}{2\pi\hbar c}.
\end{equation}

Using $L \approx 2\pi p R$ and $\mu_0 e^2 c = 4\pi\hbar\alpha$:
\begin{equation}
  \frac{U_{\text{EM}}}{E_{\text{circ}}} = \frac{4\pi\hbar\alpha \cdot R \cdot [\ln(8R/r) - 2]}{2\pi p R \cdot 2\pi\hbar} = \frac{\alpha[\ln(8R/r) - 2]}{p\pi}.
\end{equation}

For $p = 2$ and $r/R = \alpha$: $U_{\text{EM}}/E_{\text{circ}} = \alpha[\ln(8/\alpha) - 2]/(2\pi) \approx 0.0058$, or about 0.6\%. The self-energy is a perturbative correction to the dominant circulation energy.

%----------------------------------------------------------------------
\section{S3. Angular Momentum Computation}
\label{sec:S3}

\subsection{Instantaneous angular momentum}

At parameter $\lambda$, the photon has position $\mathbf{r}(\lambda)$ and momentum $\mathbf{p} = (E_{\text{total}}/c)\hat{v}(\lambda)$, where $\hat{v} = \mathbf{v}/c$ is the unit tangent. The $z$-component of angular momentum is:
\begin{equation}
  L_z(\lambda) = [\mathbf{r} \times \mathbf{p}]_z = \frac{E_{\text{total}}}{c}\left[x\hat{v}_y - y\hat{v}_x\right].
\end{equation}

\subsection{Time average}

The photon moves at constant speed $c$, so it spends equal proper time per unit arc length. The time-averaged angular momentum is:
\begin{equation}
  \langle L_z \rangle = \frac{\int_0^{2\pi} L_z(\lambda)\left|\dfrac{d\mathbf{r}}{d\lambda}\right| d\lambda}{\int_0^{2\pi}\left|\dfrac{d\mathbf{r}}{d\lambda}\right| d\lambda}.
\end{equation}

This is evaluated numerically with $N = 2000$--$4000$ sample points.

\subsection{Sweep over \texorpdfstring{$r/R$}{r/R}}

Holding $R$ fixed and sweeping $r/R$ from 0 to 1:

\begin{table}[h]
\centering
\begin{tabular}{@{}ccl@{}}
\toprule
$r/R$ & $\langle L_z \rangle / \hbar$ & Physical regime \\
\midrule
0.001 & 0.9997 & Nearly circular (spin-1) \\
0.01 & 0.997 & Thin torus \\
$\alpha \approx 0.0073$ & 0.500 & \textbf{Spin-$\frac{1}{2}$} \\
0.1 & 0.91 & Moderate torus \\
0.3 & 0.73 & Thick torus \\
0.5 & 0.57 & \\
0.9 & 0.15 & Fat torus (spin $\to$ 0) \\
\bottomrule
\end{tabular}
\caption{Angular momentum as a function of aspect ratio.}
\label{tab:angular-momentum}
\end{table}

The transition from integer spin to half-integer spin occurs at $r/R = \alpha$ --- the same value fixed by the self-energy self-consistency condition. The angular momentum constraint and the electromagnetic self-energy constraint point to the same geometry.

\subsection{Bisection for simultaneous solution}

The code performs a two-step procedure:
\begin{enumerate}
  \item \textbf{Find $r/R$ where $\langle L_z \rangle = \hbar/2$} at fixed $R$, using bisection on $r/R$.
  \item \textbf{Find $R$ where $E_{\text{total}} = m_e c^2$} at the fixed $r/R$ found in step~1, using bisection on~$R$.
\end{enumerate}

Both constraints are satisfied simultaneously.

%----------------------------------------------------------------------
\section{S4. Topology Catalog}
\label{sec:S4}

Complete survey of candidate surfaces for the null worldtube:

\begin{table}[h]
\centering
\small
\begin{tabular}{@{}lccccccccl@{}}
\toprule
Surface & Genus & Orient. & $\mathbb{R}^3$ & R1 & R2 & R3 & R4 & R5 & Status \\
\midrule
Sphere $S^2$ & 0 & yes & yes & yes & \textbf{no} & \textbf{no} & yes & yes & Fails R2, R3 \\
Torus $T^2$ & 1 & yes & yes & yes & yes & yes & yes & yes & \textbf{Model} \\
Klein bottle & 0$^*$ & no & imm. & yes & yes & yes & no$^*$ & yes & Forces $p=2$ \\
M\"obius strip & --- & no & yes & yes & yes & no & yes & yes & $= T^2$ boundary \\
Double torus & 2 & yes & yes & yes & yes & yes & yes & \textbf{no}$^*$ & Curv.\ unstable \\
Genus-$g$ & $g$ & yes & yes & yes & yes & yes & yes & \textbf{no}$^*$ & Curv.\ unstable \\
Knotted torus & 1 & yes & yes & yes & yes & yes & yes & yes & Pert.\ corrections \\
\bottomrule
\end{tabular}
\caption{Topology catalog for null worldtube candidates. $^*$Curvature energy exceeds rest mass by factor $\sim 1/(3\pi\alpha) \approx 15$ for genus~2.}
\label{tab:topology}
\end{table}

The Klein bottle is not a separate surface but provides the \emph{electromagnetic} boundary condition explaining why $p$ must be even. The M\"obius strip is the Seifert surface bounded by the $(2,1)$ torus knot, explaining fermionic statistics. Knotted tori (trefoil, figure-8, etc.)\ produce mass corrections of order $\alpha^2 m_e \sim$~keV --- fine structure, not generation-scale splitting.

%----------------------------------------------------------------------
\section{S5. Klein Bottle EM Boundary Conditions}
\label{sec:S5}

\subsection{The identification}

A torus has two independent identifications:
\begin{equation}
  \text{Torus:}\quad (x, 0) \sim (x, 2\pi r) \quad\text{and}\quad (0, y) \sim (2\pi R, y).
\end{equation}

The Klein bottle reverses the second:
\begin{equation}
  \text{Klein bottle:}\quad (x, 0) \sim (x, 2\pi r) \quad\text{and}\quad (0, y) \sim (2\pi R, 2\pi r - y).
\end{equation}

\subsection{Effect on EM fields}

After one toroidal circuit ($\theta \to \theta + 2\pi$), the poloidal coordinate reverses: $\phi \to 2\pi - \phi$. Under this identification, the electric field component transverse to the circulation direction transforms as $E_\perp \to -E_\perp$ (orientation reversal).

For a standing EM mode $\Psi(\theta, \phi) \propto e^{ip\theta}e^{iq\phi}$, single-valuedness requires:
\begin{equation}
  \Psi(\theta + 2\pi, 2\pi - \phi) = \Psi(\theta, \phi),
\end{equation}
\begin{equation}
  e^{i2\pi p}e^{-iq\phi} = e^{iq\phi} \quad \Rightarrow \quad p \in 2\mathbb{Z}.
\end{equation}

Only \textbf{even} toroidal winding numbers produce single-valued electromagnetic modes on the Klein bottle. The minimum is $p = 2$, giving the $(2,1)$ torus knot as the lightest fermion.

\subsection{Mode analysis}

\begin{table}[h]
\centering
\begin{tabular}{@{}cccc@{}}
\toprule
Mode $(p,q)$ & Klein bottle & Torus & Statistics \\
\midrule
$(1,1)$ & forbidden & allowed & boson \\
$(2,1)$ & \textbf{allowed} & allowed & fermion \\
$(1,0)$ & forbidden & allowed & trivial \\
$(3,1)$ & forbidden & allowed & fermion \\
$(4,1)$ & allowed & allowed & fermion (heavier) \\
\bottomrule
\end{tabular}
\caption{Mode compatibility with Klein bottle boundary conditions.}
\label{tab:klein-modes}
\end{table}

The Klein bottle eliminates all odd-$p$ modes. The remaining modes, with even~$p$, automatically bound M\"obius strips and are fermionic. The boson--fermion distinction is thus a consequence of electromagnetic boundary conditions.

%----------------------------------------------------------------------
\section{S6. Gauss--Bonnet Stability Argument}
\label{sec:S6}

\subsection{The theorem}

For a closed orientable surface of genus $g$:
\begin{equation}
  \iint K\, dA = 2\pi\chi = 2\pi(2 - 2g),
\end{equation}
where $K$ is the Gaussian curvature and $\chi$ the Euler characteristic.

\subsection{Curvature energy}

An electromagnetic mode on a curved surface experiences a conformal coupling correction:
\begin{equation}
  E^2 \to E^2 + (\hbar c)^2\frac{K}{6}.
\end{equation}

For a surface with total area $A = 4\pi^2 Rr$ and average curvature $\langle |K| \rangle = |2 - 2g|/(2\pi Rr)$:
\begin{equation}
  \frac{\delta E_{\text{curv}}}{E} \sim \frac{(\hbar c)^2 |K|}{6E^2} = \frac{|2 - 2g|}{6\pi \cdot (r/R)} = \frac{|2 - 2g|}{6\pi\alpha}.
\end{equation}

\subsection{Results by genus}

\begin{table}[h]
\centering
\begin{tabular}{@{}ccccc@{}}
\toprule
Genus & $\chi$ & Flat? & $\delta E / E$ & Stable? \\
\midrule
0 (sphere) & 2 & no & --- & Fails R2 \\
1 (torus) & 0 & \textbf{yes} & 0 & \textbf{Stable} \\
2 & $-2$ & no & $1/(3\pi\alpha) \approx 14.6$ & Unstable \\
3 & $-4$ & no & $2/(3\pi\alpha) \approx 29.1$ & Unstable \\
$g$ & $2-2g$ & no & $(g-1)/(3\pi\alpha)$ & Unstable \\
\bottomrule
\end{tabular}
\caption{Curvature stability by genus.}
\label{tab:genus-stability}
\end{table}

The curvature instability ratio is \textbf{independent of particle mass}: it depends only on the aspect ratio $r/R = \alpha$ and the genus. Any particle --- electron, muon, tau, quark --- on a genus-2 surface with $r/R = \alpha$ has curvature energy exceeding rest mass by a factor of $\sim$15. The torus is uniquely stable.

For a genus-2 surface to be marginally stable ($\delta E / E < 1$), it would need $r/R > 1/(3\pi) \approx 0.106$ --- fifteen times the electron's aspect ratio $\alpha \approx 0.0073$.

%----------------------------------------------------------------------
\section{S7. Retarded Field Structure}
\label{sec:S7}

The self-interaction of the circulating photon involves retarded fields --- the electromagnetic field at each point on the curve due to all other points, evaluated at the retarded time.

For each point $i$ on the curve, the retarded point $j$ satisfies:
\begin{equation}
  |\mathbf{r}_i - \mathbf{r}_j| = c(t_i - t_j),
\end{equation}
with periodic wrapping on the closed worldtube ($t$ is defined modulo $T$).

At the electron's self-consistent geometry:
\begin{itemize}
  \item Mean retarded distance: $\langle |\mathbf{r}_i - \mathbf{r}_j| \rangle \sim R$ (comparable to major radius)
  \item Mean retarded delay: $\langle t_i - t_j \rangle / T \sim 0.4$ (roughly 40\% of the circulation period)
  \item Retarded distance fluctuation: $\sigma_d / \langle d \rangle \sim 0.3$ (moderate variation around the loop)
\end{itemize}

The retarded field structure confirms that the photon's self-interaction is dominated by the toroidal-scale geometry ($R$), not the tube-scale geometry ($r$). This justifies the use of the Neumann inductance formula, which captures the dominant $O(R)$ contribution.

\subsection{Physical interpretation}

The retarded field calculation reveals that the photon ``sees'' itself approximately 40\% of a circulation period in the past. For the electron, this delay is $\Delta t \approx 0.4T \approx 3.3 \times 10^{-21}$~s. The retarded self-field provides the restoring mechanism: the photon's electromagnetic field, arriving from its own past trajectory, confines the photon to the torus. This is self-consistent because the field and the source are the same object --- the photon creates the field that confines it.

The mean retarded distance $\langle d \rangle \sim R$ means the dominant self-interaction is toroidal (between opposite sides of the torus hole), not poloidal (between nearby segments of the tube). This hierarchy --- toroidal interaction at scale $R$ dominant over poloidal interaction at scale $r$ --- is what allows the self-energy to be treated as a perturbation on the circulation energy.

%----------------------------------------------------------------------
\section{S8. Effective Potential \texorpdfstring{$V_{\text{eff}}(R)$}{V\_eff(R)}}
\label{sec:S8}

The total energy as a function of torus major radius, with $r/R = \alpha$ held fixed:
\begin{equation}
  V_{\text{eff}}(R) = E_{\text{circ}}(R) + U_{\text{EM}}(R) = \frac{2\pi\hbar c}{L(R)} + \frac{\alpha[\ln(8/\alpha) - 2]}{p\pi}\cdot\frac{2\pi\hbar c}{L(R)}.
\end{equation}

Since both terms scale as $1/L(R) \propto 1/R$, the potential is \textbf{monotonically decreasing}: $V_{\text{eff}} \propto 1/R$. There are no local minima.

This means a single isolated torus has no preferred radius --- the system wants to expand indefinitely (lowering its energy by increasing $R$). Three consequences follow:

\begin{enumerate}
  \item \textbf{Single-particle stability is topological}, not energetic. The torus topology cannot be continuously deformed away; the photon is topologically trapped. The radius $R$ is not fixed by a potential minimum but by external constraints (the Koide quantization condition for three linked tori).

  \item \textbf{The Bohr atom analogy.} The Coulomb potential $V(r) = -ke^2/r$ is also monotonic (no minimum in the radial potential without the centrifugal barrier). Bohr's quantization condition $2\pi r = n\lambda_{\text{dB}}$ selects discrete radii from the continuum. Similarly, the Koide angle $\theta_K = (6\pi + 2)/9$ selects three discrete torus radii.

  \item \textbf{Quantum corrections} --- running coupling $\alpha(\mu)$, Casimir energy, and Euler--Heisenberg nonlinear QED --- are perturbative corrections of order $\alpha$ or smaller. They do not create potential wells. The monotonic $1/R$ shape persists through all corrections computed to date.
\end{enumerate}

Self-consistent radii:

\begin{table}[h]
\centering
\begin{tabular}{@{}lcccc@{}}
\toprule
Particle & Mass (MeV) & $R$ (fm) & $r$ (fm) & $R/\lambda_C$ \\
\midrule
$e$ & 0.511 & 194.2 & 1.42 & 0.503 \\
$\mu$ & 105.66 & 0.939 & 0.00685 & 0.00243 \\
$\tau$ & 1776.9 & 0.0558 & 0.000407 & 0.000144 \\
\bottomrule
\end{tabular}
\caption{Self-consistent radii for the three charged leptons.}
\label{tab:lepton-radii}
\end{table}

\subsection{Quantum corrections to \texorpdfstring{$V_{\text{eff}}$}{V\_eff}}

Three categories of quantum corrections have been evaluated:

\paragraph{Running coupling.}
One-loop QED gives $\alpha(\mu) = \alpha / [1 - (2\alpha/3\pi)\ln(\mu/m_e)]$. At the electron's torus radius, $\mu = \hbar c/R_e \approx 1$~MeV, this gives $\Delta\alpha/\alpha \approx 0.0004\%$ --- negligible. Running coupling becomes significant only at the muon and tau radii, where it contributes corrections of order $\alpha^2$.

\paragraph{Casimir energy.}
For a massless field on a torus with periods $2\pi R$ and $2\pi r$, the Casimir energy scales as $E_{\text{Casimir}} \sim (\hbar c/R) f(r/R)$. For the thin torus ($r/R = \alpha$), this is $O(\alpha)$ relative to $E_{\text{circ}}$ --- the same order as the self-energy already computed and included.

\paragraph{Euler--Heisenberg (nonlinear QED).}
In strong electromagnetic fields, photon--photon scattering generates an effective four-photon interaction. The correction scales as $\delta E_{\text{EH}}/E \sim \alpha^2(r_{\text{class}}/R)^4 \approx 5 \times 10^{-5}$ for the electron --- entirely negligible.

All quantum corrections are at most $O(\alpha)$, confirming that the monotonic $1/R$ shape of $V_{\text{eff}}$ is robust. The stability of individual generations must therefore come from the Koide quantization condition, not from potential wells in $V_{\text{eff}}(R)$.

%----------------------------------------------------------------------
\section{S9. Comparison Tables}
\label{sec:S9}

\subsection{NWT vs.\ Bohr model}

\begin{table}[h]
\centering
\begin{tabular}{@{}lll@{}}
\toprule
Aspect & Bohr model & Null Worldtube \\
\midrule
Orbiting entity & Electron (massive) & Photon (massless) \\
Central force & Coulomb $-ke^2/r$ & Self-interaction $U_{\text{EM}}$ \\
Quantization & $2\pi r = n\lambda$ & $\theta_K = (6\pi+2)/9$ \\
Potential shape & Monotonic $1/r$ & Monotonic $1/R$ \\
Stability mechanism & Resonance & Topology + Koide \\
Ground state & $n = 1$ & Electron \\
Excited states & $n = 2, 3, \ldots$ & $\mu$, $\tau$ (Koide rotation) \\
Coupling constant & Input ($\alpha$) & Output ($r/R$) \\
Spin & Added \emph{post hoc} & Emergent (\S IV) \\
\bottomrule
\end{tabular}
\caption{Comparison of the Bohr model and the Null Worldtube model.}
\label{tab:bohr-comparison}
\end{table}

\subsection{NWT vs.\ Schr\"odinger quantum mechanics}

\begin{table}[h]
\centering
\begin{tabular}{@{}lll@{}}
\toprule
Aspect & Schr\"odinger QM & Null Worldtube \\
\midrule
Electron described by & Wavefunction $\psi(\mathbf{r}, t)$ & Null curve on torus \\
Mass & Input parameter & Emergent: $m = h/(cL)$ \\
Spin & Added via Pauli matrices & Emergent: $L_z = \hbar/2$ \\
Charge & Input parameter & Topological invariant \\
$\alpha$ & Input parameter & Geometric ratio $r/R$ \\
Statistics & Spin-statistics axiom & M\"obius surface topology \\
Self-energy & Divergent (needs renorm.) & Finite (UV cutoff at $r$) \\
Generations & Not explained & Koide angle $\theta_K$ \\
SM params & 26 inputs & 3 inputs + 3 integers \\
\bottomrule
\end{tabular}
\caption{Comparison of Schr\"odinger QM and the Null Worldtube model.}
\label{tab:schrodinger-comparison}
\end{table}

\subsection{NWT vs.\ Dirac equation}

\begin{table}[h]
\centering
\begin{tabular}{@{}lll@{}}
\toprule
Aspect & Dirac equation & Null Worldtube \\
\midrule
Spinor components & 4 (postulated) & 2 circ.\ $\times$ 2 proj. \\
Zitterbewegung & $\omega = 2mc^2/\hbar$ (formal) & $\omega = 2\pi/T_{\text{tor}}$ (literal) \\
Antiparticles & Negative-energy sea & Reversed circulation \\
$g$-factor & $g = 2$ (exact, tree-level) & $g \approx 2$ (from circ.\ geom.) \\
Anomalous moment & $a_e = \alpha/(2\pi)$ (Schwinger) & Self-energy corr.\ to $L_z$ \\
CPT symmetry & Theorem from QFT axioms & Geometric: C\,=\,rev., P\,=\,mirr., T\,=\,t-rev. \\
Spin--orbit coupling & Relativistic correction & Cross-term tor./pol.\ motion \\
\bottomrule
\end{tabular}
\caption{Comparison of the Dirac equation and the Null Worldtube model.}
\label{tab:dirac-comparison}
\end{table}

\subsection{Key numerical values}

\begin{table}[h]
\centering
\begin{tabular}{@{}llll@{}}
\toprule
Quantity & Symbol & Value & Source \\
\midrule
Fine-structure constant & $\alpha$ & $7.297 \times 10^{-3}$ & Input / geometric ratio \\
Electron mass & $m_e$ & 0.5110~MeV/$c^2$ & Self-consistency target \\
Major radius & $R_e$ & 194.2~fm & Bisection \\
Minor radius & $r_e$ & 1.42~fm & $\alpha R_e$ \\
Compton wavelength & $\lambda_C$ & 386.2~fm & $\hbar/(m_e c)$ \\
$R_e / \lambda_C$ & --- & 0.503 & $\approx 1/2$ \\
Path length & $L$ & 2442~fm & $\approx 4\pi R_e$ \\
Circulation period & $T$ & $8.1 \times 10^{-21}$~s & $L/c$ \\
Toroidal period & $T_{\text{tor}}$ & $4.1 \times 10^{-21}$~s & $T/p = T/2$ \\
Zitterbewegung period & $T_{zb}$ & $4.1 \times 10^{-21}$~s & $2\pi\hbar/(2m_e c^2)$ \\
Current & $I$ & 19.7~A & $ec/L$ \\
Neumann inductance & $\mathcal{L}$ & $1.2 \times 10^{-18}$~H & $\mu_0 R[\ln(8R/r)-2]$ \\
Logarithmic factor & $\ln(8/\alpha) - 2$ & 4.99 & Geometric constant \\
Self-energy fraction & $U_{\text{EM}}/E_{\text{total}}$ & 0.5\% & $\alpha \ln(8/\alpha)/(2\pi p)$ \\
Torus mode count & --- & 13 & Weinberg denominator \\
\bottomrule
\end{tabular}
\caption{Key numerical values for the torus electron.}
\label{tab:numerical-values}
\end{table}

%----------------------------------------------------------------------
\begin{thebibliography}{1}

\bibitem{jackson1999}
J.~D.~Jackson,
\emph{Classical Electrodynamics}, 3rd ed.\ (Wiley, 1999).

\end{thebibliography}

\end{document}
